\begin{abstract}
On-device deep learning inference is increasingly integrated into consumer smartphone applications to guarantee low latency, user privacy, and offline functionality. While cross-platform frameworks---predominantly Google Flutter---dominate modern mobile application engineering, developers face an empirical vacuum regarding whether cross-platform abstractions degrade deep learning performance compared to native implementations. This paper presents the first controlled, full-factorial benchmark comparing Flutter and Native Android implementations of identical on-device inference workloads using Google LiteRT (formerly TensorFlow Lite). Evaluating three diverse model architectures (MobileNetV3-Large, EfficientDet-Lite0, MobileBERT-QA) across single-precision floating-point (FP32) and integer quantization (INT8), three execution delegates (multi-threaded CPU with XNNPACK, GPU, and NNAPI), and eight independent sessions per configuration, we isolate and decompose latency across data preprocessing, language-boundary foreign function interface (FFI) marshalling, native kernel execution, and postprocessing. Our empirical results demonstrate that pure kernel inference latency exhibits complete parity between Flutter and Native Android. The user-perceived end-to-end latency differential is overwhelmingly driven by data ingestion and image preprocessing in pure Dart rather than the runtime inference engine. Furthermore, while Flutter incurs fixed baseline memory and cold-start startup premiums due to its embedded virtual machine, incremental memory allocated per model arena remains equivalent. We synthesize our findings into an evidence-based decision matrix, identifying actionable thresholds for when standard Flutter is completely sufficient and when native preprocessing optimizations are justified.
\end{abstract}

\begin{IEEEkeywords}
on-device inference, edge AI, Flutter, cross-platform development, LiteRT, mobile performance, benchmarking, software engineering
\end{IEEEkeywords}
